% Source: https://artofproblemsolving.com/community/c3h2632514_water_sources
% Stable identity: aops-topic-2632514-q01
\textbf{(a)} Let the \emph{perimeter} of the water region be the number of unit edges that separate a water tile from an air tile or from the outside of the board. Initially it is at most $4k$ if John places $k$ tiles. Split each stage into single conversions, one tile at a time; each converting tile still has at least two water neighbours. If a tile with $w\ge2$ water neighbours becomes water, its $w$ shared edges stop counting and at most $4-w$ new edges start counting, so the perimeter changes by at most $4-2w\le0$. Thus the perimeter never increases. A fully flooded board has perimeter $4n$, so $4k\ge4n$ and $k\ge n$.

\textbf{(b)} Put water on the diagonal tiles $(i,i)$, $1\le i\le n$. The tile $(i,j)$ with $|i-j|=d\ge1$ touches the two tiles $(i,j\mp1)$ and $(i\pm1,j)$ pointing back towards the diagonal, which have $|i-j|=d-1$. By induction on $d$, every tile with $|i-j|=d$ becomes water exactly at stage $d+1$: it has two water neighbours from stage $d$ on, but none before. The largest value of $|i-j|$ is $n-1$, so the board is entirely water at stage $\boxed{n}$ (and not before, for $n\ge2$). Together with (a), the least number of tiles is $\boxed{n}$.
